\documentclass[11pt,a4paper]{article}
\usepackage[T1]{fontenc}
\usepackage{lmodern,amsmath,amssymb,amsthm,mathtools,booktabs,microtype}
\usepackage[margin=25mm]{geometry}
\usepackage[colorlinks=true,linkcolor=blue,citecolor=blue,urlcolor=blue]{hyperref}
\setlength{\emergencystretch}{3em}
\newtheorem{theorem}{Theorem}[section]
\newtheorem{lemma}[theorem]{Lemma}
\newtheorem{proposition}[theorem]{Proposition}
\newtheorem{corollary}[theorem]{Corollary}
\theoremstyle{remark}\newtheorem{remark}[theorem]{Remark}
\newcommand{\Q}{\mathbb Q}\newcommand{\Z}{\mathbb Z}\newcommand{\F}{\mathbb F}\newcommand{\C}{\mathbb C}
\DeclareMathOperator{\N}{N}\DeclareMathOperator{\Tr}{Tr}\DeclareMathOperator{\Span}{Span}
\title{Smaller norm determinants for rational-point encoding}
\author{}\date{14 November 2026}
\begin{document}\maketitle
\begin{abstract}
We lower the degrees of four finite integer-polynomial specifications without
changing their rational arithmetic sources, their numbers of variables, or
any arithmetic hypothesis. Norms of individual denominator factors are taken
in smaller multiquadratic subalgebras, rather than repeatedly in the entire
point algebra. A separating element for a tower of quadratic extensions
avoids quartic discriminants and needs no lower bound for nonzero radicands.
The new upper bounds are $25934337637613568<2^{55}$ in 25 variables,
$177313559789699072<2^{58}$ in 27 variables, and respectively
$126442497$ and $103165953$, both less than $2^{27}$, in the two
13-variable families. These are bounds for unexpanded polynomials.
The fully specified 27-variable member retains its certified parameter
$L=20483$ and has coefficient height less than $2^{2^{71}}$.
No reduction of a global variable count is asserted.
\end{abstract}

\section{Results and the unchanged arithmetic source}
For an integer polynomial $F(X,\boldsymbol y)$, its rational input projection
is $\{x\in\Q:\exists\boldsymbol y\in\Q^m,
F(x,\boldsymbol y)=0\}$. A count of variables includes $X$.
We use the four rational sources of \cite{UnitNorm}, with all their
nonvanishing conditions. The source itself is \emph{not} changed in this
paper. In particular, the modification to the comparison scale already
made in that paper is retained, rather than made a second time.

\begin{theorem}\label{thm:main}
Under exactly the arithmetic hypotheses of the corresponding v28 source,
there are finite integer-polynomial specifications with these bounds:
\begin{center}\small
\begin{tabular}{lrrr}\toprule
Source & Variables & Degree at most & Strict bound\\\midrule
General inherited-premise source &25&25934337637613568&$2^{55}$\\
Fully specified inherited-premise source&27&177313559789699072&$2^{58}$\\
Weak rational $abc$, exponent $121/100$&13&126442497&$2^{27}$\\
Weak rational $abc$, exponent $16/15$&13&103165953&$2^{27}$\\\bottomrule
\end{tabular}\end{center}
The new polynomial has the same rational input projection as the
corresponding v28 polynomial for every allowed fixed coefficient parameter.
Each has integer content one. In particular the same sufficient, possibly
ineffective, defining parameters work. The explicit 27-variable member at
$L=20483$ has coefficient height
\[
                              H(F)<2^{2^{71}}.
\]
\end{theorem}

The general sources retain the specified elliptic group and height inputs,
the simultaneous prime-pattern and pole-realization inputs, and, where
used, the full normalized modular realization. The conditional sources
instead use the indicated weak $abc$ inequality in their contact-height
argument. Ordinary rational $abc$ implies both indicated inequalities;
we use their precise source statements, not a new assertion about the
status of $abc$. The full external arithmetic audit and independent
mathematical review remain separate from the exact finite checks here.
The stored $L=3$ and $B=3$ members are reconstruction fixtures, not newly
certified defining members. All fixed parameters are coefficients.

The algebraic implication in Theorem~\ref{thm:main} is proved below, not
inferred from the numerical examples. The complete source circuits,
normalizations, denominator factors and coefficient recipes are in
\nolinkurl{computation/few_variables/subalgebra_norms_v30}. The source
circuits are generated by the unchanged v28 generator; their canonical
SHA-256 digests are recorded in each specification.

\section{The finite algebra and its quadratic coordinates}
Let $R=\Z[x_1,\ldots,x_a]$ be the ring of retained coordinates.
For a point pair $(h_0,h_1)$ put
\[
 A=h_0^4+12h_0^2h_1^2+4h_1^4-11,\qquad
 B=4h_0^3h_1+8h_0h_1^3+6.
\]
The point algebra is the free rank-four $R$-algebra with basis $1,U,V,W$
and multiplication
\begin{equation}\label{eq:point}
\begin{gathered}
 U^2=2A+2W,\quad V^2=A-W,\quad UV=2B,\\
 UW=AU-2BV,\quad VW=BU-AV,\quad W^2=A^2-2B^2.
\end{gathered}
\end{equation}
This is the unchanged table of \cite{UnitNorm}. Let $\mathcal B$ be the
tensor product of $p$ point blocks, of rank $b=4^p$. The reduced pole
radicands $c_1,\ldots,c_s$ are the integer elements of $\mathcal B$
specified by the v28 square extractions. The final finite algebra is
\[
 \mathcal A=\mathcal B[q_1,\ldots,q_s]/(q_j^2-c_j),\qquad n=4^p2^s.
\]
Here $(p,s)=(9,2),(9,3),(3,2),(3,2)$ in the four rows of
Theorem~\ref{thm:main}. These are ranks of free modules, even when some
square classes are dependent. Such dependence does not authorize
reducing the number of geometric characters in a split finite algebra.

The v28 construction specifies a finite set $\mathcal F\subset\mathcal B$
of nonconstant syntactic factors. It covers every numerator required to
be nonzero, every denominator, every extracted square multiplier, and
all factors of the reduced pole radicands. Fixed nonzero integer factors
are immaterial. Its former scalar condition was
\begin{equation}\label{eq:oldguard}
                 \Delta=\N_{\mathcal B/R}\left(\prod_{f\in\mathcal F}f\right)\ne0.
\end{equation}
The full-unit completeness theorem in \cite{UnitNorm} is an essential
inherited result: the required integer inputs have witnesses on this
locus. Our new scalar condition will have \emph{exactly the same locus}
at every rational retained tuple.

\subsection{Splitting the point blocks over a quadratic coefficient field}
Put $K=\Q(\sqrt2)$ and, in $\mathcal B\otimes K$, define
\[
 \xi_i=\frac{U_i+\sqrt2V_i}{2},\qquad
 \eta_i=\frac{U_i-\sqrt2V_i}{2}.
\]
Then
\begin{equation}\label{eq:split}
 \xi_i^2=A_i+\sqrt2B_i,\quad \eta_i^2=A_i-\sqrt2B_i,
 \quad U_i=\xi_i+\eta_i,\quad
 V_i=(\xi_i-\eta_i)/\sqrt2,\quad W_i=\xi_i\eta_i.
\end{equation}
The square-exponent monomials $e_v$, indexed by $v\in\F_2^{2p}$,
form a basis over $K[x_1,\ldots,x_a]$. Their multiplication adds
exponent vectors modulo two and multiplies by the scalar radicands
for repeated exponents. Thus every linear subspace
$L\subseteq\F_2^{2p}$ defines a free subalgebra
\[
             \mathcal C_L=\Span_{K[\boldsymbol x]}\{e_v:v\in L\}
\]
of rank $2^{\dim L}$. Write $\sigma$ for simultaneous interchange of
the two exponent positions in every point block, together with
coefficient conjugation in $K$.

\begin{lemma}\label{lem:smallnorm}
Suppose $f\in\mathcal B$ belongs to $\mathcal C_L$ and $L$ is stable
under that interchange. Then
\[
                         g_f=\N_{\mathcal C_L/K[\boldsymbol x]}(f)
\]
belongs to $\Z[\boldsymbol x]$. At every rational tuple where the
point algebra is etale, $g_f\ne0$ if and only if $f$ is a unit of
$\mathcal B$ at that tuple.
\end{lemma}
\begin{proof}
The indicated multiplication matrix has entries in $K[\boldsymbol x]$.
Conjugation preserves its algebra and fixes $f$, so it fixes the
determinant. Hence $g_f\in\Q[\boldsymbol x]$.
The point generators are integral over $R$, as follows also from the
monic equations for $W,U,V$ in \eqref{eq:point}. Every conjugate of
$f$ is integral over $R$. Over an algebraic closure of the fraction
field the norm is a product of some of those conjugates, so it too
is integral over $R$. Since $\Z[\boldsymbol x]$ is integrally closed,
$g_f\in R$. This argument proves integrality even if a matrix in the
split basis displays denominators two.

At an etale rational specialization, the $2^{2p}$ sign characters of
the full split algebra restrict to all characters of $\mathcal C_L$.
This can be seen directly: each character of the exponent subspace
extends to the full binary vector space. After fixing any nonzero
square roots of the radicands, their sign changes give all extensions.
The norm is nonzero precisely when every restricted character takes a
nonzero value on $f$, equivalently when every full character does so.
This is the unit condition. Descent to $\Q$ does not change it.
\end{proof}

The point algebra is etale at every rational retained tuple. Indeed,
$A_i+B_i\sqrt2=h_i^4-(3-\sqrt2)^2$ cannot vanish: taking norms
would give $\N_{K/\Q}(h_i)^4=49$, impossible for a rational norm.
Its conjugate is nonzero as well. This argument imposes no positivity
or rational-square-root hypothesis on a putative solution.

\subsection{Computing the subspaces without expanding the circuits}
Associate to an arithmetic circuit node an affine subspace containing
its possible square-exponent support. Scalars have support $\{0\}$.
For the $i$th point, $U_i,V_i$ have support contained in
$e_{2i-1}+\langle e_{2i-1}+e_{2i}\rangle$, and $W_i$ has support
$\{e_{2i-1}+e_{2i}\}$. Negation leaves the affine space unchanged.
For a sum take the affine span of the constituent spaces. For a product
add their offsets and sum their direction spaces. For a positive power
$k$ of $v+L$ use $(k\bmod2)v+L$, and for a zeroth power use $\{0\}$.
These rules contain the actual support: parity records multiplication,
and cancellations can only remove terms. Take $L_f$ to be the linear
span of the resulting affine space. The initial spaces and all these
operations are stable under conjugation, so $L_f$ is stable as well.
No claim that $L_f$ is the smallest possible space is needed.

Define
\begin{equation}\label{eq:newguard}
                            G=\prod_{f\in\mathcal F}g_f,
                 \qquad g_f=\N_{\mathcal C_{L_f}}(f).
\end{equation}
Lemma~\ref{lem:smallnorm} proves $G\ne0\iff\Delta\ne0$ at every
rational retained tuple. For example, already in one point block,
\[
 \N_{\langle1,W\rangle}(1+U^2)=1+4A+8B^2,
 \qquad \N_{\mathcal B}(1+U^2)=(1+4A+8B^2)^2.
\]
The repeated power contributes to the old degree but not to its
nonzero locus. The construction removes such repetitions separately
for every specified factor, including factors involving several points.

\section{Separation in a quadratic tower}
We first give the general assertion, including nonreal square roots.

\begin{lemma}\label{lem:tower}
Let a finite etale algebra over $\Q$ be presented as a tower with
successive generators $t_1,\ldots,t_m$, where
$t_i^2=d_i$ and $d_i$ belongs to the preceding algebra. Suppose every
geometric value of every $d_i$ is nonzero. Let $M_i\ge1$ be positive
rational numbers bounding $|t_i|$ at all geometric characters.
Define recursively
\[
 \phi_{m+1}=0,\qquad
 \phi_i=\frac{t_i}{2M_i}+\frac{d_i}{4M_i^2}\phi_{i+1}.
\]
Then the geometric values of $\phi_1$ are distinct and nonzero.
Moreover $|\phi_i|\le2/3$ on every character.
\end{lemma}
\begin{proof}
The bound follows backwards from $1/2+(1/4)(2/3)=2/3$.
For nonvanishing, the leading term has magnitude $|t_i|/(2M_i)$
and the tail at most $|t_i|^2/(6M_i^2)\le|t_i|/(6M_i)$.
Thus $|\phi_i|\ge|t_i|/(3M_i)>0$.

Compare two distinct characters, and let $i$ be the first generator
where they differ. Their earlier values, hence their $d_i$ values,
coincide; their $t_i$ values are opposite. The leading difference in
$\phi_i$ has magnitude $|t_i|/M_i$, whereas the tail difference has
magnitude at most $|d_i|/(3M_i^2)\le|t_i|/(3M_i)$.
It cannot cancel. At each earlier step the difference is multiplied
by the same nonzero factor $d_j/(4M_j^2)$, so the difference in
$\phi_1$ is nonzero. This argument uses absolute values in $\C$ and
never assumes a positive radicand or a lower bound for its magnitude.
\end{proof}

We use a tower inside the original point algebra, without replacing its
integral determinant basis. In a point block write
\[
 N=A^2-2B^2,\quad M=2A^2-B^2,\quad
 A'=6A^3-9AB^2-4B^3,\quad B'=-4A^3+6AB^2+3B^3,
\]
\[
                   z=(2A+B)U-2(A+B)V.
\]
The exact identities
\begin{equation}\label{eq:towerpoint}
 W^2=N,\qquad z^2=2A'+2MW,\qquad A'^2-M^2N=2(B')^2
\end{equation}
show that $W,z$ are successive quadratic generators. At a rational tuple,
$N,M,B'$ are nonzero. For $B'$, if $B\ne0$ use Eisenstein at three
on $4T^3-6T-3$; if $B=0$ use $A\ne0$. The two values of the second
radicand have product $8(B')^2$, so neither vanishes. To check that
no character is lost, put $m_0=2A+B$, $m_1=-A-B$. Then
\[
 U=\frac{m_0z-4m_1B'/z}{M},\qquad
 V=\frac{-m_1z+2m_0B'/z}{M}.
\]
These recover $U,V$ on every character, not just the actual real one.

Order the tower as $W_i,z_i$ in point order, followed by the reduced
pole roots. On $G\ne0$ all pole radicands are units, so
Lemma~\ref{lem:tower} applies to a tower of length $m=2p+s$ and rank
$n=2^m$. The next section constructs polynomial choices of $M_i$.
Clearing the recursive denominators gives the integral element
\begin{equation}\label{eq:omega}
 \Omega=\sum_{i=1}^m 2M_it_i
       \left(\prod_{k<i}d_k\right)
       \left(\prod_{k>i}4M_k^2\right).
\end{equation}
Indeed $\Omega=(\prod_k4M_k^2)\phi_1$. The clearing factor is a
nonzero rational scalar at every rational retained tuple, so $\Omega$
is also nonzero and separating on $G\ne0$.

\section{The polynomial and its exact rational meaning}
Take multiplication matrices in the original tensor basis of
\eqref{eq:point}, followed by exponent zero or one in each pole root.
Define
\begin{equation}\label{eq:polynomial}
                    P(\boldsymbol x,Y)=\det_{\mathcal A}(I-YG M_\Omega).
\end{equation}
It is an integer polynomial. If $G=0$, it equals one identically in
$Y$, even if a pole specialization is nonreduced. If $G\ne0$, the
algebra is etale and its geometric eigenvalues $G\Omega(\chi)$ are
distinct and nonzero. A rational zero $Y$ therefore gives a rational
simple eigenvalue $1/Y$ of a rational matrix.

For completeness of the rational descent, the corresponding kernel is
one-dimensional over $\Q$. Every multiplication matrix commutes with
$G M_\Omega$ and hence acts on that line by a rational scalar. The
assignment of that scalar is a unital algebra homomorphism
$\mathcal A\to\Q$. Thus it reconstructs a rational character of the
entire finite source algebra, including each required rational square
root. Conversely, a rational character on $G\ne0$ gives the rational
zero $Y=1/(G\Omega(\chi))$. Rational zeros of an arbitrary resultant
would not suffice for this argument; separation is the decisive step.

The reconstructed character satisfies every unchanged source guard,
because $G\ne0$ is equivalent to \eqref{eq:oldguard}. The inverse pole
rescalings and square extractions are those of \cite{UnitNorm}. Thus the
rational projection onto all retained coordinates is exactly the v28
full-unit source projection. Its completeness theorem proves the
claimed input equivalence with the same coefficient parameters.
For the conditional constructions use $F=XP$ as before, adding input
zero without a new variable. Since $P(\boldsymbol x,0)=1$, both $P$
and $XP$ have integer content one. This proves the algebraic assertion
in Theorem~\ref{thm:main}.

\section{Auditable degree and coefficient bounds}
The following bounds deliberately avoid expanding the large determinants.
All logarithms of coefficient bounds in this section are to base two.
Give retained coordinates weight one, and point basis elements $U,V,W$
weights $2,2,4$. Pole roots have the weights specified below.
If a multiplication matrix represents an element of weight at most $d$,
its entry in row $u$, column $v$ has degree at most $d+w_v-w_u$.
Weights telescope in a determinant term, giving the rank times $d$.

For coefficient norms in the split basis use
$\|a+b\sqrt2\|=|a|+2|b|$, and its polynomial coefficient sum.
This is submultiplicative. Scale each square-exponent monomial by
$8^{|v|}$. Since $\|A\|_1\le28$, $\|B\|_1\le18$, and
$28+2\cdot18=64$, multiplication by $\xi_i,\eta_i$ has induced
weighted column norm at most eight. By \eqref{eq:split}, multiplication
by $U_i,V_i,W_i$ has norm at most $16,16,64$, respectively. Thus
all v28 circuit bounds, computed with these same three caps, remain
valid in the split basis and in every invariant subspace used here.

Let $d_f,\ell_f$ be the recorded weight and logarithmic operator-norm
bounds for a guard factor, and $r_f=\dim L_f$. The determinant bound
\begin{equation}\label{eq:guardbounds}
 \deg g_f\le2^{r_f}d_f,\qquad
 \log\|g_f\|_1\le2^{r_f}\ell_f
\end{equation}
follows from the weighted column-sum bound for determinants. To justify
using a coefficient norm rather than a pointwise norm, expand the
Leibniz formula, apply submultiplicativity coefficientwise, and bound
the sum by the product of the column sums. The row/column scaling
factors cancel in every permutation. Since $g_f$ is rational, its
$K$-coefficient norm is its ordinary rational coefficient norm.
Write $g=\sum_f2^{r_f}d_f$ and $h=\sum_f2^{r_f}\ell_f$.

\subsection{Polynomial root envelopes}
For the tower radicands and generators, the following quadruples
$(d_i,\ell_i,w_i,j_i)$ give respectively radicand weight, radicand
logarithmic operator norm, root weight, and root logarithmic operator
norm:
\[
             (8,11,4,6)\quad\hbox{for }W_i,\qquad
             (12,20,6,12)\quad\hbox{for }z_i.
\]
For example $\|A^2-2B^2\|_1\le1432<2^{11}$, and
\[
 2(6\cdot28^3+9\cdot28\cdot18^2+4\cdot18^3)
       +2(2\cdot28^2+18^2)64=715552<2^{20}.
\]
The root $z$ itself has norm at most
$(2\cdot28+18)16+2(28+18)16=2656<2^{12}$.
For a pole radicand, use its source bounds $(d_i,\ell_i)$ and take
$w_i=\lceil d_i/2\rceil$, $j_i=\lceil\ell_i/2\rceil+1$.
The two-by-two companion block, with the root part of the basis scaled
by $2^{\lceil\ell_i/2\rceil}$, proves this operator bound.

Set
\[
 r=1+\sum_{j=1}^a x_j^2,\qquad
 k_i=\lceil d_i/4\rceil,\quad v_i=\lceil\ell_i/2\rceil,
 \qquad M_i=2^{v_i}r^{k_i}.
\]
These are positive integer polynomials. To verify their root bounds at
an arbitrary rational tuple, rescale basis vector of weight $w$ by
$r^{w/2}$ as well as its constant coefficient-norm scale.
The entry degree inequalities then give operator norm at most
$2^{\ell_i}r^{d_i/2}$ for multiplication by $d_i$.
Its eigenvalues are the geometric radicand values. Taking square
roots gives $|t_i|\le2^{\ell_i/2}r^{d_i/4}\le M_i$.
This uses only $r\ge1$ and $|x_j|\le\sqrt r$; all signs and all
complex embeddings are included.

Let
\[
 e_i=2k_i,\qquad v_i'=v_i+k_i\lceil\log_2(a+1)\rceil.
\]
Then $\deg M_i\le e_i$ and $\log\|M_i\|_1\le v_i'$.
The $i$th summand of \eqref{eq:omega} has weight at most
\[
 o_i=e_i+w_i+\sum_{k<i}d_k+2\sum_{k>i}e_k
\]
and logarithmic operator norm at most
\[
 l_i=1+v_i'+j_i+\sum_{k<i}\ell_k+\sum_{k>i}(2+2v_k').
\]
Put $o=\max_i o_i$ and $l=\lceil\log_2m\rceil+\max_i l_i$.
The determinant \eqref{eq:polynomial} therefore satisfies
\begin{equation}\label{eq:finalbounds}
 \deg P\le n(g+o+1),\qquad \deg_Y P\le n,\qquad
                         \log\|P\|_1\le n(h+l+1).
\end{equation}
The first term in each matrix column is the identity term; adding it
accounts for the final one in the logarithmic bound. Multiplication by
$X$ adds one to total degree and changes neither coefficient norm nor
support cardinality.

\subsection{Exact evaluated majorants}
The independent checker reconstructs every factor, its parity subspace,
all tower bounds, and the following quantities from the integer circuits:
\begin{center}\scriptsize
\begin{tabular}{lrrrr}\toprule
Source&$g$&$h$&$o$&$l$\\\midrule
25&24732739297&181440737464&172420&1695856\\
27&84549104420&708535385583220&589540&4941393118\\
13, $d=6$&487255&1215402&6660&29662\\
13, $d=4$&397571&1080876&5420&25322\\\bottomrule
\end{tabular}\end{center}
For the explicit source at $L=20483$, the resulting logarithmic
coefficient bound is
\[
               1485916763799083286528<2^{71},
\]
which proves the coefficient assertion of Theorem~\ref{thm:main}.
The other recorded coefficient bounds refer only to the stored small
fixtures; no effective defining parameters are inferred from them.
The degree analysis is uniform in the allowed fixed coefficient
parameter, because it changes coefficient nodes, not any exponent of
a retained variable or the parity-support analysis.

For $V=a+1$ variables and $D=n(g+o+1)$, an additional explicit upper
bound for the number of monomials is
\[
       \binom{D+V}{V}-\binom{D-n+V-1}{V}.
\]
It counts monomials of total degree at most $D$ whose $Y$ exponent is
at most $n$. The binomial coefficient formula subtracts those divisible
by $Y^{n+1}$. It remains valid after multiplication by $X$.
No exact expansion, exact degree, or minimality claim is made.

\section{Reconstruction, tests, and limits of the certificate}
The producer emits four complete specifications and a bound table.
The separate checker uses iterative circuit traversal, low-pivot binary
elimination, and independently transcribed recurrences; it does not
import the producer's support or bound routines. It checks all source
guards and recorded denominators, exact square extraction, every norm
rank, and every degree and coefficient-norm bound. Twenty deliberate
corruptions test omissions and lowered bounds. Symbolic tests verify the
point tower and smaller-norm identities. Nine exact rank-eight examples
include negative and tiny radicands, nonreal point choices, dependent
radicands, and a zero-norm boundary, and compare rational eigenvalues
with rational characters. A separate repeated-squareclass test detects
a nonseparating sum. In addition, at one exact split finite-field tuple
for each conditional source, all 64 point characters are evaluated. Every
recorded factor is constant on the fibers of its proposed subalgebra
characters; its Fourier support and smaller/full norm identity are checked.
These finite-field fixtures test the actual source circuits, not only
the small symbolic examples.

These tests supplement the proofs; they do not replace them. The
production determinants of ranks $2^{20}$, $2^{21}$, and $2^8$ have not
been expanded. The source arithmetic assumptions and the inherited
full-unit completeness theorem are mathematical dependencies, not
conclusions of the exact finite certificate. All earlier byte-pinned
source files remain unchanged.

\begin{thebibliography}{9}
\bibitem{UnitNorm} \emph{A single norm determinant for the 25- and
13-variable constructions}, project manuscript, 14 November 2026.
\nolinkurl{papers/finite_algebra/unit_norm_20261114_v28.tex}.
\bibitem{General} \emph{A defining polynomial in twenty-five variables:
combined scalar and comparison conditions}, project source manuscript.
\nolinkurl{papers/few_variables25/few_variables25_merged_20261110_v19.tex}.
\bibitem{Explicit} \emph{An explicit defining polynomial in twenty-seven
variables}, project source manuscript.
\nolinkurl{papers/few_variables27/explicit27_20261110_v19.tex}.
\bibitem{Conditional} \emph{Thirteen Variables from Direct Elliptic Index
Comparison}, project source manuscript.
\nolinkurl{papers/conditional_few_variables/direct_input13_20261110_v22.tex}.
\end{thebibliography}
\end{document}
